\section{Model Construction}

\subsection*{IMMC Requirement 1: Conservation Priorities and Defining Protection}

\subsubsection{Identifying and Prioritizing Conservation Challenges}

Etosha National Park faces complex, multi-dimensional threats requiring systematic prioritization. We categorize challenges by species, habitat, location, and threat type:

\textbf{Species-Specific Priorities:} We employ the Analytic Hierarchy Process (AHP) to systematically prioritize species based on IUCN status, ecological role, and conservation urgency (detailed in Section 2.1). This yields quantitative priority weights $\pi_s$ for each species.

\textbf{Habitat-Specific Risks:} Waterholes (86 sites) and boundary zones face highest anthropogenic pressure, while salt pan areas (4,800 km$^2$) present natural protection but limited accessibility.

\textbf{Location-Specific Threats:} Proaching hotspots correlate with road accessibility and distance from ranger outposts. Wildfire risk clusters in grassland regions during dry seasons.

\textbf{Threat Types:} Human-driven (poaching, unauthorized entry) vs. natural (wildfire, habitat disruption) require differentiated response strategies.

\subsection{Protection Definition and Priority Evaluation Model}

The first model layer establishes a quantitative protection framework. We define protection using a multi-criteria evaluation approach combining ecological importance, threat exposure, and operational capability.

\subsubsection{Hierarchical Structure for Protection Assessment}

We employ the Analytic Hierarchy Process (AHP) to organize protection criteria into a hierarchical structure:

\begin{equation}
\mathcal{P} = \{L_1, L_2, L_3, L_4\}
\end{equation}

where:
\begin{itemize}
\item $L_1$: Goal layer—overall protection level
\item $L_2$: Criterion layer—species priority, habitat importance, threat exposure, coverage capability
\item $L_3$: Indicator layer—IUCN status, breeding grounds, poaching risk, patrol frequency
\item $L_4$: Alternative layer—spatial units (monitoring zones)
\end{itemize}

\subsubsection{AHP Weight Calculation Method}

The AHP weight vector $\mathbf{w} = (w_1, w_2, w_3, w_4)$ is calculated through the following steps:

\textbf{Step 1: Construct Pairwise Comparison Matrix}

Using the 1-9 Saaty scale, we construct matrix $A = [a_{ij}]_{n \times n}$ where $a_{ij}$ represents the relative importance of criterion $i$ compared to criterion $j$.

\begin{equation}
A = \begin{pmatrix}
1 & a_{12} & a_{13} & a_{14} \\
1/a_{12} & 1 & a_{23} & a_{24} \\
1/a_{13} & 1/a_{23} & 1 & a_{34} \\
1/a_{14} & 1/a_{24} & 1/a_{34} & 1
\end{pmatrix}
\label{eq:pairwise_matrix}
\end{equation}

\textbf{Step 2: Calculate Weight Vector}

The weight for criterion $i$ is computed using the geometric mean method:

\begin{equation}
w_i = \frac{(\prod_{j=1}^{n} a_{ij})^{1/n}}{\sum_{k=1}^{n} (\prod_{j=1}^{n} a_{kj})^{1/n}}
\label{eq:ahp_weight}
\end{equation}

\textbf{Step 3: Consistency Verification}

Consistency ratio $CR$ must satisfy $CR < 0.1$:

\begin{equation}
CR = \frac{CI}{RI} = \frac{(\lambda_{max} - n)/(n - 1)}{RI}
\label{eq:consistency}
\end{equation}

where $\lambda_{max}$ is the maximum eigenvalue of $A$, and $RI$ is the random index ($RI = 0.90$ for $n=4$).

\subsubsection{Etosha-Specific AHP Weights}

For Etosha National Park, the pairwise comparison yields:

\begin{equation}
\mathbf{w} = (w_1, w_2, w_3, w_4) = (0.50, 0.30, 0.15, 0.05)
\label{eq:etosha_weights}
\end{equation}

\noindent where:
\begin{itemize}
\item $w_1 = 0.50$: Ecological importance (species priority + habitat quality)
\item $w_2 = 0.30$: Threat exposure (poaching risk + human intrusion)
\item $w_3 = 0.15$: Operational capability (patrol accessibility + response time)
\item $w_4 = 0.05$: Seasonal variation factor
\end{itemize}

This weighting prioritizes \textbf{conservation value} (80\% combined for ecology + threat) over operational convenience.

\subsubsection{Fuzzy Comprehensive Protection Evaluation}

Protection is not binary; many regions are only partially protected. We apply fuzzy comprehensive evaluation to capture partial protection states:

\begin{equation}
P_i = \mathbf{w} \circ \mathbf{R}_i = \sum_{j=1}^{n} w_j \cdot r_{ij}
\label{eq:fuzzy_protection}
\end{equation}

where:
\begin{itemize}
\item $\mathbf{w} = (w_1, w_2, w_3, w_4)$: AHP-derived weight vector
\item $\mathbf{R}_i = (r_{i1}, r_{i2}, r_{i3}, r_{i4})^T$: membership degree vector for zone $i$
\item $\circ$: fuzzy composition operator
\end{itemize}

\textbf{Membership Function Design:}

For criterion $j$ in zone $i$, the membership degree $r_{ij}$ is calculated as:

\begin{equation}
r_{ij} = \begin{cases}
1, & \text{if } v_{ij} \geq v_{max} \\
\frac{v_{ij} - v_{min}}{v_{max} - v_{min}}, & \text{if } v_{min} < v_{ij} < v_{max} \\
0, & \text{if } v_{ij} \leq v_{min}
\end{cases}
\label{eq:membership}
\end{equation}

where $v_{ij}$ is the observed value, and $v_{max}, v_{min}$ are the maximum and minimum values across all zones.

\textbf{Protection Score Classification:}

\begin{equation}
\text{Protection Level} = \begin{cases}
\text{High}, & \text{if } P_i \geq 0.75 \\
\text{Medium}, & \text{if } 0.50 \leq P_i < 0.75 \\
\text{Low}, & \text{if } P_i < 0.50
\end{cases}
\label{eq:protection_level}
\end{equation}

\subsubsection{Species Priority Weighting}

Species priority combines IUCN status, ecological role, and flagship value:

\begin{equation}
\pi_s = \alpha_1 \cdot S_{IUCN}(s) + \alpha_2 \cdot S_{keystone}(s) + \alpha_3 \cdot S_{flagship}(s)
\label{eq:species_priority}
\end{equation}

\noindent For Etosha, using $\alpha_1 = 0.6$, $\alpha_2 = 0.3$, $\alpha_3 = 0.1$:

\begin{table}[h]
\centering
\caption{Species Priority Weights for Etosha National Park}
\begin{tabular}{lcccc}
\hline
\textbf{Species} & \textbf{IUCN Status} & \textbf{Keystone} & \textbf{Flagship} & \textbf{Priority ($\pi_s$)} \\
\hline
Black Rhino & CR (0.95) & 0.90 & 0.95 & \textbf{0.940} \\
Elephant & EN (0.85) & 0.95 & 0.90 & \textbf{0.885} \\
Lion & VU (0.75) & 0.85 & 0.80 & \textbf{0.790} \\
Cheetah & VU (0.75) & 0.70 & 0.75 & \textbf{0.735} \\
Wild Dog & EN (0.85) & 0.75 & 0.70 & \textbf{0.785} \\
Leopard & NT (0.65) & 0.70 & 0.75 & \textbf{0.690} \\
\hline
\end{tabular}
\end{table}

\noindent IUCN status values: CR=0.95, EN=0.85, VU=0.75, NT=0.65

\subsubsection{Protection Gap Measure}

The protection gap quantifies the difference between current and target protection:

\begin{equation}
G_i = \tau_{target} - P_i
\label{eq:protection_gap}
\end{equation}

where $\tau_{target} = 0.85$ is the target protection level for high-priority zones.

\textbf{Example Calculation:}

For Zone 17 (Halali region):
\begin{itemize}
\item Species indicator: $r_{17,1} = 0.92$ (black rhino + elephant presence)
\item Habitat indicator: $r_{17,2} = 0.88$ (waterhole + breeding ground)
\item Threat indicator: $r_{17,3} = 0.45$ (moderate poaching risk)
\item Coverage indicator: $r_{17,4} = 0.67$ (limited road access)
\end{itemize}

\begin{equation}
P_{17} = 0.50 \times 0.92 + 0.30 \times 0.88 + 0.15 \times 0.45 + 0.05 \times 0.67 = 0.821
\end{equation}

\begin{equation}
G_{17} = 0.85 - 0.821 = 0.029
\end{equation}

This small gap (2.9\%) indicates Zone 17 is near-target protection and requires maintenance-level resources.

\textbf{Priority Zone Ranking:}

Zones with largest protection gaps receive highest allocation priority:

\begin{equation}
\text{Priority Rank}(i) = \text{sort}(G_i, \text{descending})
\label{eq:priority_rank}
\end{equation}

\subsection{Spatial Representation and Network Model}

The second layer represents Etosha as a spatial network, incorporating roads, entrances, waterholes, and terrain.

\subsubsection{Graph-Based Spatial Representation}

We model Etosha as a weighted graph $G = (V, E, W)$ where:
\begin{itemize}
\item $V = \{v_1, v_2, ..., v_{50}\}$: 50 monitoring zones
\item $E = \{(v_i, v_j) | \text{patrol route exists}\}$: traversable connections
\item $W = \{w_{ij} | (v_i, v_j) \in E\}$: travel time/distance weights
\end{itemize}

\subsubsection{Accessibility Matrix}

Accessibility from ranger outpost $k$ to zone $i$:

\begin{equation}
A_{ik} = \begin{cases}
1, & \text{if } d(v_k, v_i) \leq R_{patrol} \\
e^{-\lambda(d(v_k, v_i) - R_{patrol})}, & \text{if } d(v_k, v_i) > R_{patrol}
\end{cases}
\label{eq:accessibility}
\end{equation}

where $d(v_k, v_i)$ is the shortest-path distance computed via Dijkstra's algorithm.

\subsubsection{Response Time Calculation}

Response time from outpost $k$ to zone $i$:

\begin{equation}
T_{response}(i, k) = T_{travel}(i, k) + T_{prepare} + T_{uncertainty}
\label{eq:response_time}
\end{equation}

where $T_{travel}$ is shortest-path travel time, $T_{prepare}$ is preparation time (15 minutes), and $T_{uncertainty}$ accounts for weather and terrain variability.

\subsection{Resource Allocation and Planning Model}

The third and central layer is a constrained integer programming model for resource deployment.

\subsubsection{Decision Variables}

\begin{align}
x_{ijk} &= \begin{cases}
1, & \text{if patrol } k \text{ assigned to zone } i \text{ in shift } t \\
0, & \text{otherwise}
\end{cases} \label{eq:decision_patrol} \\
y_{ij} &= \begin{cases}
1, & \text{if camera trap } j \text{ placed at zone } i \\
0, & \text{otherwise}
\end{cases} \label{eq:decision_camera} \\
z_{ik} &= \begin{cases}
1, & \text{if UAV } k \text{ monitors zone } i \\
0, & \text{otherwise}
\end{cases} \label{eq:decision_uav}
\end{align}

\subsubsection{Multi-Objective Integer Programming Model}

\begin{equation}
\max_{x, y, z} \left[ \mathcal{P}_{weighted}(x, y, z), -\mathcal{C}_{cost}(x, y, z), -\mathcal{G}_{gap}(x, y, z) \right]
\label{eq:multi_objective}
\end{equation}

subject to:

\noindent \textbf{Resource Constraints:}
\begin{align}
\sum_{i,k} x_{ijk} &\leq N_{rangers} \quad \forall t \label{eq:ranger_limit} \\
\sum_{i} y_{ij} &\leq N_{cameras} \label{eq:camera_limit} \\
\sum_{i} z_{ik} &\leq N_{UAVs} \label{eq:uav_limit}
\end{align}

\noindent \textbf{Coverage Constraints:}
\begin{equation}
\sum_{k} A_{ik} \cdot x_{ijk} + \beta_{camera} \cdot \sum_{j} y_{ij} + \beta_{UAV} \cdot \sum_{k} z_{ik} \geq \tau_{min} \quad \forall i \in \mathcal{L}_{high}
\label{eq:coverage_requirement}
\end{equation}

\noindent \textbf{Temporal Constraints:}
\begin{align}
\sum_{j} x_{ijk} &\leq T_{shift} \quad \forall i, k, t \label{eq:shift_limit} \\
x_{ijk} + x_{ijk'} &\leq 1 \quad \text{(no consecutive shifts)}
\end{align}

\noindent \textbf{Budget Constraint:}
\begin{equation}
c_{ranger} \sum_{i,j,k} x_{ijk} + c_{camera} \sum_{i,j} y_{ij} + c_{UAV} \sum_{i,k} z_{ik} \leq B
\label{eq:budget}
\end{equation}

\subsubsection{Priority-Weighted Objective}

\begin{equation}
\mathcal{P}_{weighted} = \frac{\sum_{i} w_i \cdot P_i(x, y, z)}{\sum_{i} w_i}
\label{eq:weighted_protection}
\end{equation}

where $w_i = \sum_{s \in S_i} \pi_s \cdot \alpha_s(i)$ is the conservation priority weight.

\subsection{Protection Gap Minimization}

The protection gap formulation:

\begin{equation}
\min_{x, y, z} \sum_{i} w_i \cdot \max(0, \tau_{target} - P_i(x, y, z))
\label{eq:gap_minimization}
\end{equation}

This formulation directly minimizes the shortfall from target protection levels.

\subsection*{IMMC Requirement 2: Resource Allocation Model}

\subsubsection{Strategic Deployment Under Resource Constraints}

Our integer programming model directly addresses the IMMC requirement to "allocate protection resources...accounting for limited resources, geographic considerations, and uncertainty, emphasizing strategic deployment rather than complete coverage."

\textbf{Key Features of Our Allocation Model:}

\begin{enumerate}
    \item \textbf{Limited Resources:} Explicit constraints on ranger availability (127 minimum), UAV fleet (45 units), camera traps (120 units), and budget (\$1.2M annually)

    \item \textbf{Geographic Considerations:} Graph-theoretic accessibility matrix $A_{ik}$ accounts for road networks (3,551 km), terrain difficulty, and response time calculations via Dijkstra's algorithm

    \item \textbf{Uncertainty Handling:} Monte Carlo simulation (10,000 trials) propagates uncertainty in threat levels, detection probabilities, and resource availability through the optimization

    \item \textbf{Strategic Deployment:} Priority-weighted objective (Eq.~\ref{eq:weighted_protection}) concentrates resources on high-priority zones (protection gap $G_i > 0.15$) rather than uniform coverage

    \item \textbf{Multi-Method Coordination:} Hybrid deployment exploits complementarity:
    \begin{itemize}
        \item Rangers (65\% to high-priority): Intervention capability
        \item UAVs (25\% high-priority, 45\% medium-priority): Rapid broad coverage
        \item Camera traps (10\% high-priority, 60\% low-priority): Cost-effective passive monitoring
    \end{itemize}
\end{enumerate}

\textbf{Strategic vs. Complete Coverage:}

Our model deliberately \textit{does not} aim for complete spatial coverage. Instead, it maximizes protection per unit resource:
\begin{itemize}
    \item Complete coverage would require 400+ rangers (infeasible)
    \item Our strategic deployment achieves 89\% protection with 127 rangers (31\% reduction)
    \item Focus on high-priority species and critical habitats yields maximum conservation impact per dollar
\end{itemize}

\textbf{Resource Types Allocated:}

Our model allocates three categories of protection resources:
\begin{enumerate}
    \item \textbf{Personnel:} 127 rangers in 3-shift rotation (42 day, 43 afternoon, 42 night)
    \item \textbf{Monitoring Technologies:}
    \begin{itemize}
        \item UAVs (45 units): Aerial surveillance, thermal imaging, real-time transmission
        \item Camera traps (120 units): Continuous monitoring, species identification
        \item Satellite monitoring: Seasonal habitat change detection
    \end{itemize}
    \item \textbf{Other Interventions:}
    \begin{itemize}
        \item Intelligence networks: Community informant systems
        \item Rapid response teams: Quick reaction to incident detection
        \item K-9 units: Tracker dogs for poacher pursuit
    \end{itemize}
\end{enumerate}
